%% ==================== 符号说明 模板 (v1.5.6, BZD 数模社 2026 规范) ====================
%%
%% 借鉴 BZD 数模社 第四章 模板: 引导语 + 三线表 (符号/含义/单位)
%%
%% BZD 原则:
%% - 三线表只保留顶线、表头分隔线、底线 (无左右竖线, 无逐行横线)
%% - "符号"列较窄居中, "含义"列最宽左对齐, "单位"列较窄居中
%% - 无量纲统一标 "—" 或 "无量纲" (全文选一种)
%% - 数量: 通常 8-15 个核心全局符号, 半页至一页
%% - min/max/sum/ln/exp 等常见运算符不入表 (除非论文反复用)
%% - 同一变量只用一个符号, 同一符号只表示一个含义
%%
%% 跑题用户必须:
%% 1. 删掉所有 % 注释 (整段 L1-L21)
%% 2. 替换【】占位符为自己的实际符号/含义/单位
%% 3. 删 下面的 \SymbolSlot{...} 行 (check 15 抓这个 inline 【】)

\section{符号说明}

为便于后续模型的建立、求解与结果分析, 现将全文反复使用的主要符号及其含义汇总如表 \ref{tab:symbols} 所示。各模型中仅局部使用的符号将在首次出现时另行说明。

\begin{table}[htbp]
  \centering
  \caption{主要符号说明}
  \label{tab:symbols}
  \begin{tabular}{p{0.15\textwidth} p{0.55\textwidth} p{0.18\textwidth}}
    \toprule
    \textbf{符号} & \textbf{含义} & \textbf{单位} \\
    \midrule
    【$x_i$】 & 【第 $i$ 个样本的观测值】 & 【—】 \\
    【$w_j$】 & 【第 $j$ 个指标的权重】 & 【—】 \\
    【$\mu$】 & 【样本均值】 & 【按变量确定】 \\
    【$\sigma$】 & 【样本标准差】 & 【按变量确定】 \\
    【$n$】 & 【样本量】 & 【—】 \\
    【$m$】 & 【指标数量】 & 【—】 \\
    【$t$】 & 【时间变量】 & 【年/月/日】 \\
    【$R^2$】 & 【决定系数, 衡量模型拟合优度】 & 【—】 \\
    【$\text{MAE}$】 & 【平均绝对误差】 & 【按变量确定】 \\
    【$\text{RMSE}$】 & 【均方根误差】 & 【按变量确定】 \\
    【$\text{AUC}$】 & 【ROC 曲线下面积】 & 【—】 \\
    【$\alpha$】 & 【显著性水平】 & 【—】 \\
    【$P$】 & 【概率】 & 【—】 \\
    \bottomrule
  \end{tabular}
\end{table}

\noindent\textbf{注:} 表中未列出的局部符号将在正文首次出现时说明。

% 不收录: 临时变量 (只在公式附近解释) / 局部符号 (该小节首次出现时说明) / 常见运算符.
% 必查: 一个符号只对应一个含义 / 同一变量只用一个符号 / 大小写和上下标一致 / 单位和正文一致.

% 跑题后必须删下面这一行 (check 15 占位符自检)
\textbf{【这里粘贴表格内容 (10-15 个核心符号), 删掉这一行】}
